\subsection{Lookup and permutation constraints as UCS constraints}\label{ss:ucs_lookup}

In this section and in \cref{ss:ucs_r1cs}, we show how several common relations and constraints---namely lookups, permutations, R1CS, CCS, and AIR---can be written as UCS constraints. For clarity, we present these reductions in isolation; however, in practice one can combine multiple constraint types within a single UCS instance.


We begin with lookup relations.
The construction follows Section~4.4 of~\cite{Zinc}, which proves the result for any integral domain; we include it here for completeness but omit a detailed proof.
We work over rings $R$ of the form $\mainpolyring{\QQ}{\bitbound}{\degreebound}$ or
$\FF_q^{< \degreebound}[X]$ for some prime power $q$, bit-size bound $\bitbound$, and degree bound
$\degreebound$.

We define the \emph{lookup relation} over $R$ as:
\[
\rellookup \;=\;
\left\{(\idx,\inp;\wit)\ \middle|\ 
\begin{aligned}
&\idx=(\witlen,\tablelen,\tab,R),\\
&\witlen, \tablelen\ge 0,\ \tab \in R^{\tablelen},\\
&\inp=\emptyset,\ \wit=\ww \in R^{\witlen},\\
& \ww_i\in\{\tab_{j}\mid j\in [\tablelen]\} \text{ for all }  i\in [\witlen].
\end{aligned}
\right\}.
\]
Let $\idx_{\mathsf{Look}}\defeq(\witlen,\tablelen,\tab,R)$ denote an index for $\rellookup$.
We now construct a UCS index $\idx_{\mathsf{UCS}}$ such that
$(\idx_{\mathsf{Look}},\emptyset;\ww)\in \rellookup$ if and only if
$(\idx_{\mathsf{UCS}},\emptyset;\wit')\in \ucs$ for an explicit transformation $\wit'$ of $\ww$.

We set $\numrows\defeq \witlen$ and $\numconstraints\defeq 1$.
Recall that the \emph{Lagrange basis polynomial} for $i\in[N]$ is
\begin{equation}\label{e: lagrange_basis}
\lagrange_{[N],i}(T) \defeq \prod_{j\neq i}\frac{T-j}{i-j}.
\end{equation}
It is well-defined over $\QQ$ or any $\FF_q$ with $q>N$.
Define
\begin{equation}\label{eq:L_lookup}
L(\rowvar,\ZZZ)\ \defeq\  
\sum_{k\in [\witlen]}\ZZZ_{k}\cdot \lagrange_{[\witlen],k}(\rowvar)
\end{equation}
so that $L(k,\zz)=\zz_{k}$ for all $k\in [\witlen]$.
Define
\begin{equation}
\label{eq:Q_lookup}
Q_{\mathsf{Look}}(\rowvar,\ZZZ)
\ \defeq\
\prod_{j\in[\tablelen]}\big(L(\rowvar,\ZZZ)-\tab_{j}\big).
\end{equation}
%
If $R=\mainpolyring{\QQ}{\bitbound}{\degreebound}$, define
$\idx_{\mathsf{UCS}}=(\witlen,\inplen{=}0,\numrows,\primetuple{=}\emptyset,\bitbound,\degbounds{=}\degreebound,\constraints)$ with
$\constraints=(Q_{0,1},\primeideal_{0,1}):=(Q_{\mathsf{Look}},(0))$.
Here the subscript $(0,1)$ denotes prime-index $\primeindex=0$ (the $\QQ[X]$-part) and constraint-index $\constraintindex=1$ (the single constraint).

If $R=\FF_q^{<\degreebound}[X]$, define
$\idx_{\mathsf{UCS}}=(\witlen,\inplen{=}0,\numrows,\primetuple{=}q,\bitbound,\degbounds{=}(0,\degreebound),\constraints)$
with $\constraints=((Q_{0,1},\primeideal_{0,1}),(Q_{1,1},\primeideal_{1,1}))$ where
$Q_{0,1}\equiv 0$, $\primeideal_{0,1}=(0)\subseteq\QQ[X]$, $\primeideal_{1,1}=(0)\subseteq\FF_q[X]$, and
$Q_{1,1}(\rowvar,\inpvar,\ZZZ_1,\ZZZ_2)\defeq Q_{\mathsf{Look}}(\rowvar,\ZZZ_2)$.
\begin{lemma}[Lookup as UCS]
\label{lem:lookup_as_ucs}
Let $\idx_{\mathsf{Look}}$ and $\idx_{\mathsf{UCS}}$ be defined as above.
For every $\ww\in R^{\witlen}$, letting $\inp=\emptyset$, we have:
\begin{itemize}[leftmargin=*]
\item If $R=\mainpolyring{\QQ}{\bitbound}{\degreebound}$ then
$(\idx_{\mathsf{Look}},\emptyset;\ww)\in\rellookup$ if and only if
$(\idx_{\mathsf{UCS}},\emptyset;\ww)\in\ucs.$
\item If $R=\FF_q^{<\degreebound}[X]$ then, writing $\mathbf{0}$ for the all-zero $\QQ$-witness
component of length $\witlen$,
$(\idx_{\mathsf{Look}},\emptyset;\ww)\in\rellookup$ if and only if
$(\idx_{\mathsf{UCS}},\emptyset;(\mathbf{0},\ww))\in\ucs.$
\end{itemize}
\end{lemma}
The proof is straightforward and we omit it. The key observation is that $R$ is an integral domain (both $\mainpolyring{\QQ}{\bitbound}{\degreebound}$ and $\FF_q^{<\degreebound}[X]$ are polynomial rings over fields, hence have no zero divisors).
Thus $Q_{\mathsf{Look}}(i,\ww)=\prod_{j}(\ww_i - \tab_j)=0$ if and only if $\ww_i\in\{\tab_j\}$. 

\begin{remark}[Degree blowup and practical lookup arguments]
The polynomial $Q_{\mathsf{Look}}$ in \eqref{eq:Q_lookup} has degree $\tablelen$ in the selector
value $L(\rowvar,\ZZZ)$.
The purpose of the above construction is to exhibit lookups as UCS instances, so that generic UCS
compilers and soundness statements apply.
In protocols, we will not prove \eqref{eq:Q_lookup} via a high-degree polynomial identity;
instead, we will use lookup arguments that exploit the lookup semantics directly, following the
approach of~\cite{Zinc}. 
\end{remark}

\begin{remark}[Lookup constraints on polynomial expressions of witness entries]
\label{rem:lookup_algebraic_combinations}
The lookup relation $\rellookup$ defined above requires each entry of a witness vector $\ww$ to belong to a table $\tab$.
In applications, one sometimes needs to constrain \emph{polynomial expressions} of witness entries to belong to a lookup table.
Concretely, given a witness vector $\ww\in R^{\witlen}$ and a polynomial $P\in R[\ZZZ]$, one may require that 
\[
  P(\ww_i)
  \;\in\; \{\tab_j \mid j\in[\tablelen]\}
  \quad\text{for all } i\in[\witlen].
\]
The extension to this setting is straightforward: one replaces the selector polynomial $L(\rowvar,\ZZZ)$ in~\eqref{eq:L_lookup} with
\[
  L'(\rowvar,\ZZZ)
  \ \defeq\
  P\Bigl(\sum_{i\in[\witlen]} \ZZZ_i \cdot \lagrange_{[\witlen],i}(\rowvar)\Bigr),
\]
and defines $Q'_{\mathsf{Look}}$ analogously.
This generalization is needed in the definition of the relation UAIR$^+$ and in some applications, notably in the SHA-256 arithmetization.
\end{remark}

\begin{remark}[Well-definedness as a lookup constraint]
\label{rem:well_defined_as_lookup}
\label{sss:well_defined_condition}
Recall that \cref{def:ucs} includes a well-definedness condition~\eqref{eq:ucs_welldefinedness} requiring $\ff_0\in(\mainpolyring{\local{q_1\cdots q_{\numprimes}}}{\bitbound}{\degreebound_0})^{\witlen}$.
This is simply a lookup constraint: the table is $\tab\defeq\mainpolyring{\local{q_1\cdots q_{\numprimes}}}{\bitbound}{\degreebound_0}$, and each entry of $\ff_0$ must belong to this finite set. Thus, the well-definedness condition can be enforced via a $\QQ[X]$-constraint in UCS and does not require separate syntax in the definition of UCS.
\begin{definition}[Well-formed UCS index]
\label{def:ucs_well_formed_index}
A UCS index $\idx$ is \emph{well-formed} if it includes a $\QQ[X]$-constraint that enforces the well-definedness condition~\eqref{eq:ucs_welldefinedness}.
\end{definition}

In many applications, other lookup constraints placed on $\ff_0$ already subsume the well-definedness condition, so it does not need to be explicitly included.
Note that while $\tab$ may be very large (or even infinite in the limit), in practice the table size is bounded by the constraint parameters.
\end{remark}




\begin{remark}[Permutation constraints as UCS instances]
\label{rem:perm_as_ucs}
So-called \emph{copy constraints} of the form $\uu_i = \vv_{\sigma(i)}$, where the permutation $\sigma:[n]\to[n]$ is fixed by the index, can also be encoded as UCS instances.
Such constraints capture copy relations that arise from wiring in arithmetizations such as Plonk~\cite{GWC19}.
We define the \emph{permutation relation} over a ring $R$ as
\[
\relperm
\;=\;
\left\{(\idx,\inp;\wit)\ \middle|\ 
\begin{aligned}
&\idx=(n,\sigma,R),\\
&\inp=\emptyset,\quad \wit=(\uu,\vv)\in R^{n}\times R^{n},\\
&\uu_i=\vv_{\sigma(i)}\text{ for all }i\in[n].
\end{aligned}
\right\}.
\]
To write this relation as a UCS instance, we view the witness as a length-$2n$ vector $\ww = (\uu || \vv)$ and define the selector polynomial
\[
Q_{\mathsf{Perm}}(\rowvar,\ZZZ)\ \defeq\
\sum_{i\in[n]}
\bigl(\ZZZ_{i}-\ZZZ_{n+\sigma(i)}\bigr)\cdot \lagrange_{[n],i}(\rowvar).
\]
Then evaluating at $\rowvar = i\in [n]$ yields $Q_{\mathsf{Perm}}(i,\ww) = \uu_i - \vv_{\sigma(i)}$.
Thus, $Q_{\mathsf{Perm}}(i,\ww) = 0$ for all $i \in [n]$ if and only if $(\uu,\vv) \in \relperm$.
The UCS construction then proceeds exactly as in \cref{lem:lookup_as_ucs}: for $R = \mainpolyring{\QQ}{\bitbound}{\degreebound}$, one sets $\constraints = (Q_{\mathsf{Perm}}, (0))$; for $R = \FF_q^{<\degreebound}[X]$, one places the constraint in the finite-field component with $Q_{1,1}(\rowvar,\inpvar,\ZZZ_1,\ZZZ_2) \defeq Q_{\mathsf{Perm}}(\rowvar,\ZZZ_2)$.
\end{remark}

\subsection{R1CS, CCS, and AIR as UCS instances}\label{ss:ucs_r1cs}

In this section we explain how several standard relations---R1CS, CCS, AIR (over fields, rationals, or polynomial rings)---can be written as $\ucs$ relations. 
Since the rings $\mainpolyring{\QQ}{\bitbound}{\degreebound}$ and $\FF_q^{<\degreebound}[X]$ are integral domains, the constructions and proofs follow directly from~\cite[Section~4.4]{Zinc}, which establishes the result for arbitrary integral domains; we summarize the key ideas below.

All transformations below work in a \emph{strictly algebraic regime}: for every constraint we set $\primeideal_{i\constraintindex}=(0)$, so that a constraint $Q_{i\constraintindex}(\cdots)\in\primeideal_{i\constraintindex}$ reduces to the polynomial identity $Q_{i\constraintindex}(\cdots)=0$.

\parhead{R1CS}
Recall that the R1CS relation over a ring $R$ requires $(M_A\cdot\zz)\circ(M_B\cdot\zz)-(M_C\cdot\zz)=\mathbf{0}$, where $M_A,M_B,M_C\in R^{\numrows\times(\inplen+\witlen)}$ are constraint matrices and $\zz=(\yy||\ww)$ concatenates the public input and witness.
The transformation to UCS proceeds as follows: first, use bivariate Lagrange interpolation to encode each matrix $M\in\{M_A,M_B,M_C\}$ as a polynomial $\hat M(\rowvar,Z)$; then define a linear selector $L_M(\rowvar;\inpvar,\ZZZ_1)$ that evaluates to the $\rowindex$-th entry of $M\cdot\zz$ when $\rowvar=\rowindex$; finally, set the UCS constraint polynomial to
\[
Q_{\ronecs}(\rowvar,\inpvar,\ZZZ_1)
\;\defeq\;
L_{M_A}\cdot L_{M_B} - L_{M_C}.
\]
For R1CS over $\mainpolyring{\QQ}{\bitbound}{\degreebound}$, the UCS index has $\primetuple=\emptyset$ and a single $\QQ[X]$-constraint $(Q_{\ronecs},(0))$.
For R1CS over $\FF_q^{<\degreebound}[X]$, one instead uses an $\FF_q[X]$-constraint with a trivial $\QQ[X]$-part.
The equivalence $(\idx_{\mathsf{R1CS}},\inp;\wit)\in\relronecs \Leftrightarrow (\idx_{\mathsf{UCS}},\inp;\wit')\in\ucs$ follows from the defining property of Lagrange interpolation; see~\cite[Section~4.4]{Zinc} for the full proof.

\parhead{CCS and AIR}
CCS relations~\cite{ccs} admit an analogous transformation as R1CS relations. Further, 
since AIR relations~\cite{ethstark} can be written as CCS relations~\cite{ccs}, they too can be instantiated as UCS instances.
An AIR-like instantiation tailored to the UCS framework, called UAIR$^+$, is described in the Zinc paper~\cite{Zinc}.

